%! AP Statistics — Unit 3, Lesson 3.3 — Lesson Plan
\documentclass[10pt]{article}
\usepackage{apstats-article}
\usepackage{apstats-boxes}

\newcommand{\UnitNumberName}{Unit 3: Collecting Data \quad}
\newcommand{\LessonNumberName}{Lesson 3.3: Random Sampling and Data Collection}

\begin{document}
\begin{center}
    {\Large\bfseries \CourseName: \SchoolYear\ (\MeetingLength) \\
    {\small\bfseries \UnitNumberName \LessonNumberName}}
\end{center}

\begin{tcolorbox}[colback=sky, colframe=navy, boxrule=0.9pt, arc=2mm,
    left=2mm, right=2mm, top=1.5mm, bottom=1.5mm]
    \textbf{Primary Objective:} Students will identify and describe the major random
    sampling methods --- SRS, stratified, cluster, systematic, and census --- and
    explain why a particular method is or is not appropriate for a given study
    (Big Idea: DAT).
\end{tcolorbox}

\renewcommand{\arraystretch}{1.6}
\begin{skillbox}[Priority Ideas \& Skills]{goldbox}
\begin{minipage}[t]{0.48\linewidth}
    \textbf{AP Skill 1 --- Selecting Statistical Methods}
    \begin{itemize}
        \item Identify a sampling method given a study description (DAT-2.C, Skill 1.C).
        \item Explain why a method is or is not appropriate (DAT-2.D, Skill 1.C).
        \item Distinguish sampling with and without replacement (DAT-2.C.1).
    \end{itemize}
\end{minipage}
\hfill
\begin{minipage}[t]{0.48\linewidth}
    \textbf{Key Understandings}
    \begin{itemize}
        \item Simple random sample: every group of size $n$ has equal chance of selection.
        \item Stratified: population divided into homogeneous strata; SRS within each.
        \item Cluster: population divided into heterogeneous clusters; SRS of clusters,
              then all individuals in selected clusters.
        \item Systematic: every $k$th individual after a random start.
        \item Census: every individual in the population is measured.
    \end{itemize}
\end{minipage}
\end{skillbox}

\begin{skillbox}[{Vocabulary, Concepts, \& Theorems}]{greenbox}
\begin{minipage}[t]{0.48\linewidth}
    \begin{itemize}
        \item \textbf{Simple random sample (SRS)} --- every group of size $n$ has an
              equal chance of being chosen
        \item \textbf{Stratified random sample} --- population divided into strata
              (homogeneous groups); SRS within each stratum
        \item \textbf{Cluster sample} --- population divided into heterogeneous clusters;
              a random sample of clusters is selected; all individuals in chosen
              clusters are surveyed
    \end{itemize}
\end{minipage}
\hfill
\begin{minipage}[t]{0.48\linewidth}
    \begin{itemize}
        \item \textbf{Systematic random sample} --- random start, then every $k$th
              individual
        \item \textbf{Census} --- all items/subjects in a population are measured
        \item \textbf{Sampling with/without replacement} --- whether a selected
              individual can be chosen again
        \item \textbf{Sampling frame} --- the list from which the sample is drawn
    \end{itemize}
\end{minipage}
\end{skillbox}

\begin{skillbox}[Activate Prior Knowledge \& Spiral Review]{sky}
\begin{minipage}[t]{0.48\linewidth}
    \textbf{Prior Knowledge Check (3--4 min)}
    \begin{itemize}
        \item Recall: ``What's the difference between an observational study and an
              experiment?'' (Lesson 3.2)
        \item Recall: ``Why does random selection make data more trustworthy?''
              (Lesson 3.1)
        \item Transition: ``Today we get specific about HOW to randomly select
              individuals.''
    \end{itemize}
\end{minipage}
\hfill
\begin{minipage}[t]{0.48\linewidth}
    \textbf{Connections Forward}
    \begin{itemize}
        \item When random sampling fails or isn't used $\to$ Lesson 3.4 (bias)
        \item Random \emph{assignment} in experiments $\to$ Lesson 3.5
        \item Scope of inference (generalization) $\to$ Lesson 3.7
        \item Inference for proportions (Unit 7) uses SRS assumptions
    \end{itemize}
\end{minipage}
\end{skillbox}

\begin{skillbox}[Hook \& Lesson]{sky}
\begin{minipage}[t]{0.48\linewidth}
    \textbf{Hook (5--7 min)}\\
    Challenge: ``I need to survey 50 of the 400 students in this school. How would
    you do it?'' Collect 3--4 student ideas (pull names from a hat, ask friends,
    survey every 8th person in the hallway, etc.). Classify each as random or not.\\[4pt]
    Transition: \textit{``Statisticians have developed several principled methods --- all
    relying on chance --- that work better than asking your friends.''}
\end{minipage}
\hfill
\begin{minipage}[t]{0.48\linewidth}
    \textbf{Lesson Flow (15--18 min)}
    \begin{enumerate}
        \item SRS: define; show numbering + random number table or calculator.
        \item Stratified: define strata (homogeneous); SRS within each; advantage ---
              guarantees representation of each group.
        \item Cluster: define cluster (heterogeneous); random selection of clusters;
              advantage --- practical when individuals are spread out geographically.
        \item Systematic: random start, every $k$th individual; caution about periodic
              patterns in the list.
        \item Census: everyone; when practical, eliminates sampling variability.
        \item Pros/cons table comparing all methods.
    \end{enumerate}
\end{minipage}
\end{skillbox}

\begin{skillbox}[Explicit Instruction \& Active Monitoring]{sky}
\begin{minipage}[t]{0.48\linewidth}
    \textbf{Direct Instruction (10 min)}
    \begin{itemize}
        \item Worked example: a school of 600 students (100 per grade 9--12, plus
              100 each in two elective tracks). Choose 60 students using SRS, stratified
              by grade, cluster by homeroom, and systematic.
        \item Emphasize: stratified $\neq$ cluster --- strata are homogeneous (same
              grade), clusters are heterogeneous (mixed grades/backgrounds).
    \end{itemize}
\end{minipage}
\hfill
\begin{minipage}[t]{0.48\linewidth}
    \textbf{Active Monitoring}
    \begin{itemize}
        \item Listen for: confusing stratified and cluster (most common error).
        \item Cold-call: ``If I want to compare grades 9 and 11 specifically, which
              method ensures I have students from each grade?'' (stratified)
        \item Check that students can describe the procedure, not just name the method.
    \end{itemize}
\end{minipage}
\end{skillbox}

\begin{skillbox}[Group Work \& Differentiation]{redbox}
\begin{minipage}[t]{0.48\linewidth}
    \textbf{Partner Activity (8--10 min)}\\
    Four scenario cards (hospital, city block, school, factory). Partners:
    \begin{enumerate}
        \item Name the sampling method described.
        \item Explain one advantage of using that method in this context.
        \item Identify one disadvantage or situation where a different method would
              be better.
    \end{enumerate}
\end{minipage}
\hfill
\begin{minipage}[t]{0.48\linewidth}
    \textbf{Differentiation}
    \begin{itemize}
        \item \textit{Support}: Comparison table with columns: Method / Key feature /
              When to use / One advantage.
        \item \textit{Extension}: Design two different sampling plans for a citywide
              health survey of 1{,}000 residents. Justify which is better and why.
        \item \textit{ELL}: Visual diagrams of each method (color-coded circles
              representing strata/clusters).
    \end{itemize}
\end{minipage}
\end{skillbox}

\begin{skillbox}[Individual Work \& Assessment]{redbox}
\begin{minipage}[t]{0.48\linewidth}
    \textbf{Exit Ticket (5 min)}\\
    Scenario: \textit{``A researcher wants to survey 80 of 500 registered voters in
    a county. She divides them into 5 districts (100 per district) and randomly
    selects 2 districts. She then surveys all 100 voters in each selected
    district.''}
    \begin{enumerate}
        \item Name the sampling method.
        \item What is one advantage of this method?
        \item What is one disadvantage?
    \end{enumerate}
\end{minipage}
\hfill
\begin{minipage}[t]{0.48\linewidth}
    \textbf{AP-Style Check}\\
    \textit{``A magazine surveys its subscribers by sending a questionnaire to every
    25th name on the subscriber list, starting from a randomly chosen name between 1
    and 25. This is an example of a:''}
    \begin{enumerate}[label=(\Alph*)]
        \item Simple random sample
        \item Stratified random sample
        \item Cluster sample
        \item Systematic random sample
    \end{enumerate}
    Answer: (D)
\end{minipage}
\end{skillbox}

\begin{skillbox}[Reinforcement \& Extension]{goldbox}
\begin{minipage}[t]{0.48\linewidth}
    \textbf{Homework / Practice}
    \begin{itemize}
        \item For 4 sampling scenarios: identify the method, justify whether it is
              appropriate, and state what population results can be generalized to.
        \item Comparison question: ``In what situation would stratified random sampling
              be better than SRS?''
    \end{itemize}
\end{minipage}
\hfill
\begin{minipage}[t]{0.48\linewidth}
    \textbf{Extension / Preview}
    \begin{itemize}
        \item \textit{Extension}: Design a sampling plan for a national political poll.
              Explain which method you would use and why.
        \item \textit{Preview}: Lesson 3.4 examines what goes wrong when random
              sampling is not used or is poorly executed --- start thinking about
              whose voices get missed.
    \end{itemize}
\end{minipage}
\end{skillbox}

\end{document}
